% Shared title metadata and result phrases for talk.tex and script.tex.
% Every magnitude below is sourced from progress_brief.tex (P2's results
% write-up); background/model phrasing is sourced from research_brief.tex.
% Numbers are never rounded further for the slide than they are in the brief.

\newcommand{\PaperTitle}{Network Position and Academic Performance\\[0.15em]
  {\normalsize Adding Network Centrality to Smirnov \& Thurner (2017)}}
\newcommand{\PaperShortTitle}{Network Position and Academic Performance}
\newcommand{\PaperAuthors}{Amy Bing \and Tingying Huang}
\newcommand{\AuthorShort}{Bing \& Huang}
\newcommand{\PresenterName}{Tingying Huang}
\newcommand{\CoauthorsSpoken}{Amy Bing}
\newcommand{\PaperAffiliations}{Monash University}
% Shared affiliation for both authors -- the simple default path (no
% per-author \AuthorAffil tabular needed since both are Monash).
\newcommand{\TitleAuthorBlock}{\PaperAuthors}
\newcommand{\TitleAffiliationBlock}{\PaperAffiliations}
\newcommand{\KeyQuestion}{Does network position predict grades?}
\newcommand{\KeyQuestionSubtitle}{On top of who a student's direct friends are}
\newcommand{\VenueDate}{ECC3841 Network Economics}
\newcommand{\VenueShort}{Presentation 2 --- Progress Report}
\newcommand{\TitleDisclaimer}{}

% ---- Result phrases: slide (written) and spoken (script) forms share the ----
% ---- same source-traced number, worded for the medium.                  ----

% Result 1 -- replication. Source: progress_brief.tex, "Result 1".
\newcommand{\ResultOneSlide}{Own past GPA strongly predicts next GPA ($\beta=0.955$); a friend's past GPA does not (\(p=0.557\))}
\newcommand{\ResultOneSpoken}{a student's own past GPA is a very strong predictor of their next GPA, but once you control for that, a friend's past GPA predicts nothing}

% Result 2 -- naive effect is popularity. Source: progress_brief.tex, "Result 2".
\newcommand{\ResultTwoSlide}{Katz-Bonacich alone: $+0.093$ ($p<0.001$). Add degree: $-0.028$ ($p=0.446$)}
\newcommand{\ResultTwoSpoken}{on its own, centrality looks like it matters, but the moment you add raw popularity to the same regression, that effect disappears}

% Result 3 -- pooled headline, read at face value. Source: progress_brief.tex, "Result 3".
\newcommand{\ResultThreeSlide}{$\beta=-0.0415$, SE\(=0.019\), \(p=0.033\), \(n=36{,}696\)}
\newcommand{\ResultThreeSpoken}{a small, negative, statistically significant effect, at the value of phi we fixed in advance}

% Result 4 -- corrected test. Source: progress_brief.tex, "Result 4".
\newcommand{\ResultFourSlide}{No \(\phi\) is significant ($p=0.36$ to $0.93$); the sign is not even consistent}
\newcommand{\ResultFourSpoken}{at every value of phi, not significant, and the sign flips back and forth}

% Headline / conclusion. Source: progress_brief.tex, boxed conclusion.
\newcommand{\ResultHeadlineSlide}{Net of popularity, friend GPA, and a student's own past performance, we find no evidence that network position independently predicts grades}
\newcommand{\ResultHeadlineSpoken}{once we properly control for popularity, friend GPA, and a student's own past performance, there is no evidence that network position predicts grades on its own}
\newcommand{\ResultSupportSlide}{Naive centrality looked positive, the pooled sweep looked negative, the corrected test looks like nothing}
\newcommand{\ResultSupportSpoken}{the naive test looked positive, the pooled test looked negative, and the one correctly specified test looks like nothing at all}
\newcommand{\ResultImplicationSlide}{Interventions built around widening students' network reach are not supported by this data once the obvious confounds are handled}
\newcommand{\ResultImplicationSpoken}{so policies that try to help students by widening their network reach are not supported by what we've found here}
